% !TEX program = pdflatex
% !BIB program = bibtex
\documentclass[11pt]{article}

% Change "review" to "final" to generate the final (sometimes called camera-ready) version.
% Change to "preprint" to generate a non-anonymous version with page numbers.
\usepackage[review]{acl}

% Standard package includes
\usepackage{times}
\usepackage{latexsym}

% For proper rendering and hyphenation of words containing Latin characters (including in bib files)
\usepackage[T1]{fontenc}
% For Vietnamese characters
% \usepackage[T5]{fontenc}
% See https://www.latex-project.org/help/documentation/encguide.pdf for other character sets

% This assumes your files are encoded as UTF8
\usepackage[utf8]{inputenc}

% This is not strictly necessary, and may be commented out,
% but it will improve the layout of the manuscript,
% and will typically save some space.
\usepackage{microtype}

% This is also not strictly necessary, and may be commented out.
% However, it will improve the aesthetics of text in
% the typewriter font.
\usepackage{inconsolata}

%Including images in your LaTeX document requires adding
%additional package(s)
\usepackage{graphicx}
\usepackage{amsmath}
\usepackage{amssymb}
\usepackage{booktabs}
\usepackage{multirow}
\usepackage{float}
\usepackage{placeins}
\usepackage{stfloats}
\usepackage{adjustbox}
\usepackage[nameinlink,noabbrev]{cleveref}
\emergencystretch=2em
\hfuzz=20pt
\vfuzz=20pt
\hbadness=10000
\vbadness=10000
% If the title and author information does not fit in the area allocated, uncomment the following
%
%\setlength\titlebox{<dim>}
%
% and set <dim> to something 5cm or larger.

\usepackage{soul}
\usepackage{color, xcolor}
\definecolor{lightblue}{RGB}{0, 191, 255}
\sethlcolor{lightblue}
\soulregister\cite7
\soulregister\eqref7
\soulregister\ref7 
\newcommand{\mathcolorbox}[2]{\colorbox{#1}{$\displaystyle #2$}}
\newcommand{\clr}[1]{\textcolor{blue}{#1}}

\title{A Multi-cluster Boundary Learning Method for Out-of-Scope Intent Detection via MiniLM Embedding}

% \author{
%   \textbf{Yihong Xu\textsuperscript{1}},
%   \textbf{Mingyu Kang\textsuperscript{1*}},
%   \textbf{Linyuan L\"u\textsuperscript{1*}}\\
%   \textsuperscript{1}University of Science and Technology of China, Hefei, China\\
%   \texttt{xuyihong@mail.ustc.edu.cn}\\
%   \texttt{kangmingyu@ustc.edu.cn}\\
%   \texttt{linyuan.lv@ustc.edu.cn}
% }
% Author information can be set in various styles:
% For several authors from the same institution:
% \author{Author 1 \and ... \and Author n \\
%         Address line \\ ... \\ Address line}
% if the names do not fit well on one line use
%         Author 1 \\ {\bf Author 2} \\ ... \\ {\bf Author n} \\
% For authors from different institutions:
% \author{Author 1 \\ Address line \\  ... \\ Address line
%         \And  ... \And
%         Author n \\ Address line \\ ... \\ Address line}
% To start a separate ``row'' of authors use \AND, as in
% \author{Author 1 \\ Address line \\  ... \\ Address line
%         \AND
%         Author 2 \\ Address line \\ ... \\ Address line \And
%         Author 3 \\ Address line \\ ... \\ Address line}

% \author{First Author \\
%   Affiliation / Address line 1 \\
%   Affiliation / Address line 2 \\
%   Affiliation / Address line 3 \\
%   \texttt{email@domain} \\\And
%   Second Author \\
%   Affiliation / Address line 1 \\
%   Affiliation / Address line 2 \\
%   Affiliation / Address line 3 \\
%   \texttt{email@domain} \\}



%\author{
%  \textbf{First Author\textsuperscript{1}},
%  \textbf{Second Author\textsuperscript{1,2}},
%  \textbf{Third T. Author\textsuperscript{1}},
%  \textbf{Fourth Author\textsuperscript{1}},
%\\
%  \textbf{Fifth Author\textsuperscript{1,2}},
%  \textbf{Sixth Author\textsuperscript{1}},
%  \textbf{Seventh Author\textsuperscript{1}},
%  \textbf{Eighth Author \textsuperscript{1,2,3,4}},
%\\
%  \textbf{Ninth Author\textsuperscript{1}},
%  \textbf{Tenth Author\textsuperscript{1}},
%  \textbf{Eleventh E. Author\textsuperscript{1,2,3,4,5}},
%  \textbf{Twelfth Author\textsuperscript{1}},
%\\
%  \textbf{Thirteenth Author\textsuperscript{3}},
%  \textbf{Fourteenth F. Author\textsuperscript{2,4}},
%  \textbf{Fifteenth Author\textsuperscript{1}},
%  \textbf{Sixteenth Author\textsuperscript{1}},
%\\
%  \textbf{Seventeenth S. Author\textsuperscript{4,5}},
%  \textbf{Eighteenth Author\textsuperscript{3,4}},
%  \textbf{Nineteenth N. Author\textsuperscript{2,5}},
%  \textbf{Twentieth Author\textsuperscript{1}}
%\\
%\\
%  \textsuperscript{1}Affiliation 1,
%  \textsuperscript{2}Affiliation 2,
%  \textsuperscript{3}Affiliation 3,
%  \textsuperscript{4}Affiliation 4,
%  \textsuperscript{5}Affiliation 5
%\\
%  \small{
%    \textbf{Correspondence:} \href{mailto:email@domain}{email@domain}
%  }
%}

\begin{document}
\maketitle

\begin{abstract}
Intent detection is a critical task that bridges human intents and system actions in human-machine interaction systems. However, there still exist challenges for detecting out-of-scope (OOS) intents. (i) The traditional methods view the OOS intent detection as a multi-class classification, then the detection accuracy decreases as the class number of the known intents increases; (ii) LLM-embedding methods require large parameters, that makes them difficult to train and practically deploy. Thus, this work proposes a multi-cluster boundary learning method to detect OOS intents via MiniLM embedding (i.e., \texttt{all-MiniLM-L6-v2}) in an one-class classification workflow. The method learns the boundaries of multi-cluster embeddings generated by MiniLM from the training utterances, and then rejects the out-of-domain utterances as OOS intents. Experiments are conducted on public CLINC150, StackOverflow and Banking77 datasets. The results show that the method achieves the state-of-the-art OOS intent detection performance compared the other baselines. Ablation studies are also conducted and the results show that the used MiniLM can better adapt to the workflow and utterance embedding requirements. \clr{The code is available at supplementary materials.}
\end{abstract}

\section{Introduction}
Intent detection is a critical yet challenging task for human-machine interaction systems. It maps user utterances to system actions, which bridges human intents and system behaviors~\citep{tur2010what,louvan2020recent,wolflein2025agents}. However, human intents are complex and cannot be fully enumerated in a closed-set intent detection system. That means if the system receives an out-of-scope (OOS) intent utterance but fails to reject it, that will trigger incorrect actions and induce harmful consequences, as shown in figure~\ref{fig:motivation}. Thus, a more important and difficult task is to detect the OOS intents and reject them in a timely manner, which is actually an open intent detection task~\citep{hoffman2024inferring, muzahid2024survey}.

\begin{figure}[t]
    \centering
    \includegraphics[width=\columnwidth]{figures/bg 02.png}
    \caption{A diagram of OOS intent detection}
    \label{fig:motivation}  
\end{figure}

There are three types of intent detection methods, including statistical methods, deep-learning methods and LLM-embedding methods. Among them, the statistical methods formulate intent detection as a supervised text classification problem. These methods combine manually designed lexical features with traditional statistical classifiers, such as support vector machine~\citep{haffner2003optimizing} and naive Bayes classifier~\citep{schuurmans2020intent}. But the deep-learning methods use deep neural networks, e.g., convolutional neural network (CNN)~\citep{kim2014convolutional}, recurrent neural network (RNN)~\citep{liu2016attention} and transformer network~\citep{vaswani2017attention}, to perform intent detection via semantic feature extraction. Early deep-learning classifiers are trained on a closed-set mode, and then identify unknown out-of-domain utterances as known intents~\citep{hendrycks2017baseline, guo2017calibration}. Thus, some types of open intent detection are further proposed to conduct safe rejection for the unknown utterances with OOS intents. They are class-wise open classification method~\citep{shu2017doc}, distance-based scoring method~\citep{lee2018simple}, contrastive representation learning method~\citep{zeng2021modeling} and adaptive boundary learning method~\citep{li2025multi}. After that, with the strong semantic representation ability of LLMs, the LLM-embedding methods are proposed to extract utterance features by semantic embedding, prompt-based reasoning and uncertainty-aware agent routing~\citep{arora2024intent}.

However, the existing methods still face two limitations. First, the OOS queries with similar semantic meaning as known intents are easily absorbed into the decision regions of the known intents~\citep{li2025multi}, but in fact they are OOS. Thus, if mix OOS labels into known intent labels, and simply view the OOS intent detection as a multi-class classification, the accuracy of OOS rejection would decrease as the class number increases. Second, although LLM-embedding methods improve the ability of semantic understanding, they still require very large parameters, that induces remarkable computational costs and high sensitivity of prompt designs~\citep{zaera2025efficient}. Thus, they are difficult to train and deploy in real-time resource-constrained dialogue systems.

To address these limitations, this work proposes a multi-cluster boundary learning method for OOS intent detection with a cascade workflow. The workflow decouples the complex open intent detection problem into multiple simpler stages. This allows the workflow to use a lightweight MiniLM instead of LLM wih large-scale parameters. Moreover, it reduces the OOS intent detection to a one-class classification, instead of a complex multi-class classification. Moreover, we found that MiniLM can embed the observation data into multiple clusters. Thus, if the boundaries of clusters are learned, the out-of-domain utterances can be identified as OOS intents.

Thus, the main contributions are as follows:
\begin{enumerate}
    \item[1.] A multi-cluster boundary learning method is proposed. This method learns the boundaries of multi-cluster embeddings generated by MiniLM from the training utterances, and then rejects the out-of-domain utterances as OOS intents. 
    
    \item[2.] A cascade workflow is proposed to first conduct OOS intent detection as an one-class classification, and then conduct known intent detection as a close-set multi-class classification. This means the OOS intent detection task can be addressed independently. 

    \item[3] Experiments are conducted on public real-world datasets for OOS intent detection task. The results show that the method achieves the state-of-the-art performance for OOS intent detection. Moreover, ablation studies are also conducted, and the results show that the used MiniLM, \texttt{all-MiniLM-L6-v2}, can better adapt to the workflow and utterance embedding requirements.  
\end{enumerate}

\section{Related Work}
The open intent detection methods can be mainly categorized into three types, i.e., statistical methods, deep-learning methods and LLM-embedding methods, according to the problem settings and the types of classifiers.

\subsection{Statistical Methods for Intent Detection}
\label{sec: cascade}
The statistical methods use hand-crafted lexical features with statistical classifiers to detect the intent classes. Among them, \citet{wang2002combination} proposes a method combining rule-based grammars with statistical classifiers, but the grammar must be hand-crafted for each domain. To reduce the dependence, \citet{haffner2003optimizing} use support vector machine as discriminative model, which shows strong performance on intent classification without rule-based components. But these methods still rely on manually designed features. Thus, \citet{tur2010what} proposes an n-gram method to model local contextual dependencies, but feature sparsity limits its effectiveness on diverse intent expressions. Then, \citet{schuurmans2020intent} compares these methods, and shows that dense word embeddings outperform the sparse features for intent classification.

\subsection{Deep-Learning Methods for Intent Detection}
Then, the deep-learning methods use artificial neural networks to parametrize intent classifiers and learn dense semantic representations. E.g., CNN is used to extract the local semantic patterns from word embeddings~\citep{kim2014convolutional}, and RNN is used to capture the sequential dependencies for joint intent detection and slot filling~\citep{liu2016attention}. Moreover, Transformer model is used to directly capture the global semantic dependencies through self-attention mechanism~\citep{vaswani2017attention}. These methods improve the in-distribution accuracy, but the classification mode is still closed-set, and the classifier can only identify unknown utterances with known intent labels, and cannot reject the OOS intents~\citep{hendrycks2017baseline, guo2017calibration}. 

Thus, to reject the OOS intents, there are several methods, that view the OOS intent detection as a open-set classification task. E.g., \citet{hendrycks2017baseline} set an extra OOS class, and transform the open-set classification into a closed-set one. Then, the method assigns regular scores to each class through softmax probability. Also, \citet{lee2018simple} propose a Mahalanobis-based scoring method, that models the score as class-conditional Gaussian distribution. But these score assumptions usually violate the real-world class distributions. Moreover, with the increasing class number, the OOS labels couple with many known intent labels, then the score of OOS class is diluted and the accuracy decreases. To improve the representation separability, \citet{zeng2021modeling} propose a supervised contrastive learning, that learns discriminative embeddings to better separate intent categories. But this still yields limited improvements. 

Then, \citet{shu2017doc} drop the transformation and propose a open classification method. The method replaces the softmax probability with one-vs-rest sigmoid classifiers, and then calculates a class-wise rejection threshold. If the softmax score is over the threshold, the input utterance would be rejected. Then, \citet{zhang2021adaptive} propose an decision boundary learning method to learn a boundary from all embedding known intents, and the external region is OOS. But the method can only learn a single centroid for each intent. Actually, the embeddings usually have multiple centroids based on our observation. Then, \citet{li2025multi} propose a multi-granularity boundary learning method to learning multiple boundaries for each intent.


\subsection{LLM-Embedding Methods for Intent Detection}
Then, with the LLM technique, the LLM-embedding methods are proposed to strengthen the semantic representations and feature embeddings for user utterance by multi-agent routing mode. Among them, \citet{arora2024intent} proposes a hybrid method, that uses a Sentence Transformer to route uncertain intent predictions to a designated LLM. But this LLM is only designed for uncertain queries and is not trained for OOS queries. Then, \citet{zaera2025efficient} propose an uncertainty-aware routing method, that only triggers a fine-tuned LLM according to the Monte Carlo dropout criterion. But the routing decisions are still unstable if there exists distribution drift. And, the multi-agent LLMs require large computation cost, that make it difficult to deploy for real-world applications. Moreover, \citet{chen2025collaborating} propose a small-large model collaboration method for few-shot intent detection. The method uses LLM to generate OOS utterance and conduct data augmentation during pretraining. And then, it uses the smaller deep-learning models to realize intent detection. 

Thus, the existing methods still face two problems. (i) the OOS intent detection is usually viewed as a multi-class classification and hardly separated from the known intent detection. (ii) The multi-agent LLMs require remarkably large computation cost. And due to the reason (i), the current routing mechanism is not accurate. 


\section{Method}
\subsection{Cascade Workflow for Intent Detection}
The proposed method uses a three-stage cascade: gate $\rightarrow$ router $\rightarrow$ expert, as shown in figure~\ref{fig:model_architecture}. Given an input utterance $x$, the gate module first determines the OOS rejection. If $x$ is rejected, the gate module outputs an OOS label, otherwise, $x$ is inputted into the following modules. Then, the router module discriminates the coarse-grained class domain for the remained known intents. Then the utterance $x$ is delivered to an expert module w.r.t the class domain for fine-grained discrimination of a few classes. Note that, due to the fact that the known intent classification is close-set, all intents for training is known, thus the router and experts are all parametrized by SmolLM-135M, that is supervisedly fine-tuned through low-rank adaptation (LoRA)~\citep{hu2022lora} technique on the observation data of known intents.

\begin{figure*}[t]
    \centering
    \includegraphics[width=\textwidth]{figures/overview 01.eps}
    \caption{An Overview of the cascaded workflow via MiniLM-based multi-cluster boundary learning.}
    \label{fig:model_architecture}
\end{figure*}

For an input utterance $x$, the gate $G(\cdot)$ first checks whether $x$ is OOS:
\begin{equation}
z = G(x), \quad z \in \{\mathrm{ID}, \mathrm{OOS}\},
\end{equation}
where $z$ is the binary decision, $\mathrm{ID}$ means the input is in-distribution, and $\mathrm{OOS}$ means the input is out-of-scope. If $z = \mathrm{OOS}$, an OOS label is returned. Otherwise, the router $R(\cdot)$ predicts the domain:
\begin{equation}
d = R(x), \quad d \in \mathcal{D},
\end{equation}
where $\mathcal{D}$ is the domain set. Then the expert $E_d(\cdot)$ classifies the intent:
\begin{equation}
y = E_d(x), \quad y \in \mathcal{Y}_d \subseteq \mathcal{Y}_{\mathrm{ID}},
\end{equation}
where $\mathcal{Y}_{\mathrm{ID}}$ is the known intent label set and $\mathcal{Y}_d$ is the intent subset for domain $d$.

The proposed cascade workflow separates OOS detection from known-intent classification. Each module handles a simpler sub-task to reduces the complexity. As a result, the Gate only needs to determine whether an input belongs to the known-intent space. Thus, a lightweight model is sufficient.

\subsection{MiniLM-Based Intent Embedding}
Since the gate only needs to distinguish OOS from known-intent samples, a lightweight MiniLM encoder, \texttt{all-MiniLM-L6-v2}, is used to construct the semantic representation space. The MiniLM maps each input utterance $x$ to an embedding with fixed dimension:
\begin{equation}
e = f_{\theta}(x),\quad e \in \mathbb{R}^{h},
\label{eq: embed}
\end{equation}
where $f_{\theta}(\cdot)$ denotes the MiniLM-based encoder parameterized by $\theta$, and $h$ denotes the embedding dimension. Note that, the MiniLM with parameter $\theta$ has no need to train in this work, that implicitly assumes the MiniLM can directly embed the utterance into multiple clusters, and in fact, that is indeed the case through empirical experiments. Then, the gate applies $L_2$ normalization before computing distances to the centroids:
\begin{equation}
\bar{e} = \frac{e}{\|e\|_2}.
\end{equation}
The normalized embedding $\bar{e}$ is then used to construct the known-intent geometry for OOS detection.
    
The embedding space is not used solely as a generic feature space. It also defines the local structure of known intents.But a single-centroid assumption can be too restrictive.Because utterances with the same intent label may differ in wording, expression, and semantic focus. Thus, forcing a known intent into a single compact region can produce an inaccurate boundary. This increases the risk of rejecting valid in-domain samples or accepting OOS samples. Thus, each known intent is represented by multiple local clusters, and each cluster has a centroid.

For intent $y$, its centroid set is defined as:
\begin{equation}
\mathcal{C}_{y}=\{c_{y,1}, c_{y,2}, \ldots, c_{y,K_y}\},
\label{eq: ky}
\end{equation}
where $c_{y,k}, k=1, \dots, K_y$, denotes the $k$-th centroid of known intent $y$, and $K_y$ is the number of centroids. The centroids are constructed via K-means clustering within each intent. The number of centroids is controlled by a global setting and can be overridden for specific intents. If the number of known intents is $M$ and each intent has $K$ centroids, the gate maintains $M \times K$ centroids in total. These centroids are local geometric representatives of the known intents. Thus the representation provides the basis for boundary learning.

\subsection{Multi-cluster Boundary Learning for OOS Intent Detection}
The multi-cluster representation defines local semantic regions for known intents, and OOS detection is therefore formulated as a boundary learning problem in the embedding space. An input is accepted if it falls inside at least one learned known-intent cluster region and is rejected if it lies outside all known-intent regions. To achieve this, the gate learns a local boundary around the centroid of each cluster to separate supported inputs from OOS inputs.

For the centroid of each cluster $c_{y,k}$, the gate estimates a local radius $r_{y,k}$ from the training samples assigned to that specific cluster. In the default configuration, the radius is computed as follows:
\begin{equation}
r_{y,k}=\mu_{y,k} + \lambda \sigma_{y,k},
\label{eq: lambda}
\end{equation}
where $\mu_{y,k}$ and $\sigma_{y,k}$ are the mean and standard deviation of the distances from the assigned samples to $c_{y,k}$, and $\lambda$ controls the boundary width.

The distance function is a diagonal Mahalanobis distance. For an embedding $\bar{e}$ and a centroid $c$, the distance is computed as:
\begin{equation}
d(\bar{e},c)=\sqrt{\sum_j(\bar{e}_j-c_j)^2\frac{1}{\sigma_j^2+\epsilon}},
\end{equation}
where $\bar{e}_j$ is the $j$-th component of $\bar{e}$ and $c_j$ is the same componenet of $c$. Moreover, $\sigma_j^2$ is the corresponding feature-wise variance estimated from samples assigned to the corresponding cluster, and $\epsilon$ is a small regularization constant. This distance weights each embedding dimension by its estimated variance, which reduces the influence of dimensions with large natural variation.

For an incoming utterance, the gate compares its normalized embedding with all known-intent centroids and computes the normalized nearest-boundary score:
\begin{equation}
s(\bar{e})=\min_{y \in \mathcal{Y}_{\mathrm{ID}},\; 1 \le k \le K_y}\frac{d(\bar{e},c_{y,k})}{r_{y,k}},
\label{eq: ryk}
\end{equation}
where $c_{y,k}$ denotes the $k$-th local centroid of known intent $y$, and $r_{y,k}$ denotes the corresponding local boundary radius. The ratio of Eq.~(\ref{eq: ryk}) measures the distance from $\bar{e}$ to this centroid relative to the learned local boundary size.

The gate decision is:
\begin{equation}
G(x)=
\begin{cases}
\mathrm{ID}, & s(\bar{e}) \le 1 \\
\mathrm{OOS}, & s(\bar{e}) > 1
\end{cases}
\end{equation}
The threshold is one because each distance is normalized by its radius. If $s(\bar{e}) \le 1$, the input falls inside at least one known-intent region and is accepted. If $s(\bar{e}) > 1$, the input lies outside all known-intent regions and is rejected as OOS. The gate only decides whether an input belongs to the known intent space. It does not assign a specific intent label.

\section{Experiments}
\subsection{Datasets}
Three datasets are used for performance evaluation, as shown in Table~\ref{tab:dataset_stats}. For each dataset, they are split into train, validate and test sets under 6:1:3.

\textbf{CLINC150} is a large-scale multi-domain dataset with 150 in-scope intents and explicit OOS samples~\citep{larson2019evaluation}. The data source is available at \url{https://github.com/clinc/oos-eval}.

\textbf{StackOverflow} is a technical-domain short-text dataset with 20 intent classes and substantial semantic overlap across classes. The data source is available at \url{https://github.com/jacoxu/StackOverflow}.

\textbf{Banking77} is a fine-grained, single-domain dataset with 77 intent classes. It focuses on high-granularity classification in the financial domain. The data source is available at \url{https://github.com/PolyAI-LDN/task-specific-datasets}.

\begin{table}[H]
\centering
\renewcommand{\arraystretch}{1.05}
\resizebox{\columnwidth}{!}{
\begin{tabular}{lccc}
\toprule
\textbf{Dataset} & \textbf{\#Intents} & \textbf{\#Samples} & \textbf{Type} \\
\midrule
CLINC150       & 150              & 22500 & Multiple domains \\
StackOverflow  & 20                     & 20000 & Technical \\
Banking77  & 77 & 13083 & Banking \\
\bottomrule
\end{tabular}
}
\caption{Dataset statistics.}
\label{tab:dataset_stats}
\end{table}

\subsection{Baselines and Experimental Settings}

This work defines Known Intent Ratio (KIR) to describe the class number of all intents that split into the known intent class from datasets, as follows:
\begin{equation}
\mathrm{KIR}=\frac{|\mathcal{Y}_{\mathrm{ID}}|}{|\mathcal{Y}|},
\end{equation}
where $\mathcal{Y}$ is the full intent label set and $\mathcal{Y}_{\mathrm{ID}}$ is the subset of intents treated as known during training.

Moreover, the following state-of-the-art methods are selected as baselines: 

\textbf{MSP} \citep{hendrycks2017baseline} uses the maximum softmax probability as the confidence score for detecting misclassified or out-of-distribution samples.

\textbf{OpenMax} \citep{bendale2016towards} recalibrates activation scores and estimates the probability that an input belongs to an unknown class.

\textbf{DOC} \citep{shu2017doc} replaces the softmax layer with independent sigmoid classifiers and performs class-wise open-set rejection.

\textbf{DeepUnk} \citep{lin2019deep} uses margin loss to learn discriminative intent features and applies density-based novelty detection for unknown intent detection.

\textbf{KNNCL} \citep{zhou2022knn} combines KNN-based contrastive representation learning with density-based outlier detection for OOD intent classification.

\textbf{ADB} \citep{zhang2021adaptive} learns adaptive class-specific decision boundaries for open intent detection.

\textbf{DA-ADB} \citep{zhang2023learning} improves ADB with distance-aware representation learning and adaptive boundary learning.

Moreover, for the proposed cascade workflow, the model in gate module is 22M-parameter \texttt{all-MiniLM-L6-v2}. The models in router and expert modules are both 135M-parameter \texttt{SmolLM-135M}. The rank of LoRA (see section~\ref{sec: cascade}) is set as 32 for router, and set as 16 for expert. The scaling factor of LoRA is set as 64 for router, and set as 32 for expert. Due to the usage of LoRA, the parameter scaling is about 160M, which is close to that of BERT-base~\citep{zeng2021modeling, zhang2021adaptive, li2025multi}. The models are all trained on NVIDIA RTX 5070 GPU. The optimizer is AdamW with default settings. The batch size is 32. The learning rate is $2\times 10^{-4}$. The sub-centroid number for each intent $y$ in Eq.~(\ref{eq: ky}) is set as $K_y=2$. The boundary scaling factor in Eq.~(\ref{eq: lambda}) is set as $\lambda = 0.5$ for CLINC150 and $\lambda = 1$ for the other datasets. The embedding dimension in Eq.~(\ref{eq: embed}) is set as $h=384$.


\subsection{Evaluation Metrics}

To evaluate and compare the performances of all methods, Known F1 score, OOS F1 score and Acc are selected as metrics. 

\textbf{Known F1 score}~\citep{wang2021texsmart, zawbaa2024improved} is calculated by averaging F1 scores over all known intent classes, as follows:
\begin{equation}
\mathrm{Known\ F1}=\frac{1}{|\mathcal{Y}_{\mathrm{ID}}|}\sum_{y \in \mathcal{Y}_{\mathrm{ID}}}\frac{2P_yR_y}{P_y+R_y},
\end{equation}
where $\mathcal{Y}_{\mathrm{ID}}$ denotes the known-intent label set. $P_y$ and $R_y$ denote the precision and recall of class $y$, respectively.

\textbf{OOS F1 score}~\citep{wang2021texsmart, zawbaa2024improved} is calculated for evaluating OOS detection performance, as follows:
\begin{equation}
\mathrm{OOS\ F1}=\frac{2 P_{\mathrm{OOS}} R_{\mathrm{OOS}}}{P_{\mathrm{OOS}} + R_{\mathrm{OOS}}},
\end{equation}
where
\begin{equation}
\begin{aligned}
P_{\mathrm{OOS}}&=\frac{\mathrm{TP}_{\mathrm{OOS}}} {\mathrm{TP}_{\mathrm{OOS}}+\mathrm{FP}_{\mathrm{OOS}}}, \\
R_{\mathrm{OOS}}&=\frac{\mathrm{TP}_{\mathrm{OOS}}}
{\mathrm{TP}_{\mathrm{OOS}}+\mathrm{FN}_{\mathrm{OOS}}}.
\end{aligned}
\end{equation}
$\mathrm{TP}_{\mathrm{OOS}}$, $\mathrm{FP}_{\mathrm{OOS}}$ and $\mathrm{FN}_{\mathrm{OOS}}$ denote true positives, false positives and false negatives w.r.t. the OOS class, respectively.

\textbf{Acc} is the overall accuracy over all test samples, as follows:
\begin{equation}
\mathrm{Acc}=\frac{1}{N}\sum_{i=1}^{N}\mathbb{I}(\hat{y}_i = y_i),
\end{equation}
where $N$ is the size of test set, $y_i$ and $\hat{y}_i$ are the ground truth and the predicted labels of the $i$-th sample, respectively. And $\mathbb{I}(\cdot)$ is the indicator function.

\subsection{Performance on OOS Intent Detection}

The performance on OOS intent detection is presented in table~\ref{tab:main_results_all}. The results show that the proposed method achieves the state-of-the-art OOS F1 scores across all settings. Compared to the baselines, the proposed method achieves 0.85\%$\sim17.12\%$ improvements on OOS F1 scores.

As shown in table~\ref{tab:main_results_all}, the OOS F1 scores of all methods decrease as KIR increases, that means the dense known intent distributions intensify the semantic overlap between in-domain and OOS queries. But the proposed method maintains relatively stable OOS F1 scores under different KIR settings. That means localized cluster-wise rejection can preserve clearer semantic boundaries under dense intent distributions. Moreover, the performance gains are particularly significant on Banking77, which indicates that the localized sub-centroids can model fine-grained semantic boundaries more precisely. However, the strict front-end rejection slightly decreases Known F1, because some marginal in-domain samples are rejected as OOS queries.

\begin{table*}[!t]
\centering
\tiny
\setlength{\tabcolsep}{2.8pt}
\renewcommand{\arraystretch}{0.94}
\resizebox{\textwidth}{!}{
\begin{tabular}{c|l|ccc|ccc|ccc}
\toprule
\multirow{2}{*}{\textbf{KIR}}
& \multirow{2}{*}{\textbf{Method}}
& \multicolumn{3}{c|}{\textbf{CLINC150}}
& \multicolumn{3}{c|}{\textbf{StackOverflow}}
& \multicolumn{3}{c}{\textbf{Banking77}} \\
\cmidrule{3-11}
& & \textbf{Known F1} & \textbf{OOS F1} & \textbf{Acc}
& \textbf{Known F1} & \textbf{OOS F1} & \textbf{Acc}
& \textbf{Known F1} & \textbf{OOS F1} & \textbf{Acc} \\
\midrule
\multirow{8}{*}{$0.25$}
& MSP      & 51.02 & 59.26 & 53.38 & 42.66 & 11.66 & 27.91 & 50.47 & 39.42 & 42.19 \\
& OpenMax  & 62.65 & 77.51 & 70.27 & 47.51 & 34.52 & 38.97 & 53.42 & 48.52 & 47.76 \\
& DOC      & 75.46 & 90.78 & 86.08 & 56.30 & 62.50 & 57.75 & 65.16 & 76.64 & 70.31 \\
& DeepUnk  & 76.95 & 91.61 & 87.18 & 47.39 & 36.87 & 40.03 & 64.97 & 76.98 & 70.68 \\
& KNNCL    & 78.85 & 93.56 & 89.87 & 41.79 & 15.26 & 28.65 & 65.54 & 79.34 & 73.01 \\
& ADB      & 77.85 & 92.36 & 88.30 & 77.62 & 90.96 & 86.75 & 70.92 & 85.05 & 79.33 \\
& DA-ADB   & \textbf{79.57} & 93.20 & 89.48 & \textbf{80.87} & 92.65 & 89.07 & 73.05 & 86.57 & 81.19 \\
& \textbf{Ours} & 71.75 & \textbf{95.01} & \textbf{90.45} & 73.61 & \textbf{94.47} & \textbf{91.04} & \textbf{75.83} & \textbf{93.99} & \textbf{89.07} \\
\midrule
\multirow{8}{*}{$0.50$}
& MSP      & 72.82 & 63.71 & 66.68 & 66.28 & 26.94 & 53.23 & 73.20 & 46.29 & 61.67 \\
& OpenMax  & 79.83 & 82.15 & 80.22 & 69.88 & 46.11 & 60.27 & 75.16 & 55.03 & 65.53 \\
& DOC      & 83.84 & 87.45 & 85.19 & 77.37 & 71.18 & 73.88 & 78.38 & 72.66 & 74.60 \\
& DeepUnk  & 83.30 & 87.48 & 84.95 & 67.67 & 35.80 & 55.46 & 75.61 & 67.80 & 71.01 \\
& KNNCL    & 83.25 & 87.85 & 85.32 & 61.50 & 8.50 & 45.38 & 75.16 & 67.21 & 70.41 \\
& ADB      & 85.12 & 88.60 & 86.54 & 85.32 & 87.70 & 86.49 & 81.39 & 79.43 & 79.61 \\
& DA-ADB   & \textbf{85.58} & 90.10 & \textbf{87.93} & \textbf{86.71} & 88.86 & \textbf{87.78} & \textbf{82.54} & 79.93 & \textbf{81.51} \\
& \textbf{Ours} & 79.95 & \textbf{91.96} & 86.78 & 75.48 & \textbf{89.71} & 85.54 & 74.90 & \textbf{88.23} & 78.98 \\
\midrule
\multirow{8}{*}{$0.75$}
& MSP      & 83.65 & 63.86 & 76.19 & 81.42 & 37.86 & 73.20 & 84.99 & 46.05 & 77.08 \\
& OpenMax  & 71.14 & 75.18 & 75.36 & 82.98 & 49.69 & 75.78 & 85.50 & 53.02 & 78.32 \\
& DOC      & 87.91 & 83.87 & 85.93 & 85.64 & 65.32 & 80.55 & 84.14 & 63.51 & 78.94 \\
& DeepUnk  & 86.57 & 82.67 & 84.61 & 80.51 & 34.38 & 71.56 & 81.65 & 50.57 & 74.73 \\
& KNNCL    & 86.14 & 82.05 & 84.12 & 76.16 & 7.19 & 65.01 & 81.76 & 51.42 & 74.78 \\
& ADB      & \textbf{88.97} & 84.85 & 86.99 & 86.91 & 74.10 & 82.89 & \textbf{86.44} & 67.34 & \textbf{81.39} \\
& DA-ADB   & 88.43 & 86.00 & \textbf{87.39} & \textbf{87.66} & 74.55 & 83.56 & 85.93 & 69.37 & 81.12 \\
& \textbf{Ours} & 66.32 & \textbf{87.10} & 79.83 & 80.18 & \textbf{75.57} & 81.32 & 70.28 & \textbf{86.49} & 79.83 \\
\bottomrule
\end{tabular}
}
\caption{Main results under different known intent ratios.}
\label{tab:main_results_all}
\end{table*}

\subsection{Analysis for Gate Stage}

\begin{figure}[t]
    \centering
    \includegraphics[width=0.9\linewidth]{figures/gate_embedding_umap_real_radius.eps}
    \caption{An example of MiniLM embeddings from CLINC150 dataset when $\mathrm{KIR}=0.50$.}
    \label{fig:gate_embedding}
\end{figure}

\begin{figure}[t]
    \centering
    \includegraphics[width=\columnwidth]{figures/gate_score_distribution_clean.eps}
    \caption{Gate score distribution on CLINC150 under $\mathrm{KIR}=0.50$. The dashed line marks $s(\bar{e})=1$.}
    \label{fig:gate_score_distribution}
\end{figure}

As shown in figure~\ref{fig:gate_embedding}, MiniLM embeddings form multiple local clusters are within the same intent category. That means intent embeddings are locally concentrated rather than globally compact. Thus, the hard OOS queries may remain close to known intent regions, making single boundary modeling insufficient under complex semantic overlap. To address this problem, the proposed method maps each intent category with multiple local clusters. As shown in figure~\ref{fig:gate_score_distribution}, most known intent samples fall inside the learned cluster boundaries, while most of OOS samples fall outside the regions. That means the localized geometric modeling separates the known-intent and OOS regions effectively. However, the overlap near the threshold mainly comes from ambiguous utterances around cluster edges. That means the gate decisions are strongly influenced by local semantic structures.

\subsection{Ablation studies on MiniLM selection for Gate stage}

Ablation studies are also conducted with $\mathrm{KIR} \in \{0.25, 0.50, 0.75\}$. Three variants are evaluated: \textbf{Without Gate} replaces front-end rejection with validation-tuned confidence thresholding. \textbf{Cascade-MiniLM} uses MiniLM-based modules in all stages. \textbf{Cascade-SmolLM} uses SmolLM-based modules in all stages.

\begin{table*}[t]
  \centering
  \small
  \setlength{\tabcolsep}{0pt}
  \renewcommand{\arraystretch}{0.9}
  \begin{tabular*}{\textwidth}{@{\extracolsep{\fill}}clcccccc@{}}
  \toprule
  \multirow{2}{*}{\textbf{KIR}}
  & \multirow{2}{*}{\textbf{Variant}}
  & \multicolumn{2}{c}{\textbf{CLINC150}}
  & \multicolumn{2}{c}{\textbf{StackOverflow}}
  & \multicolumn{2}{c}{\textbf{Banking77}} \\
  \cmidrule(lr){3-4}
  \cmidrule(lr){5-6}
  \cmidrule(lr){7-8}
  & & \textbf{Acc} & \textbf{OOS F1}
  & \textbf{Acc} & \textbf{OOS F1}
  & \textbf{Acc} & \textbf{OOS F1} \\
  \midrule
  \multirow{4}{*}{$0.25$}
  & Full Pipeline  & \textbf{90.45} & \textbf{95.01} & \textbf{91.04} & \textbf{94.47} & \textbf{89.07} & \textbf{93.99} \\
  & Without Gate   & 67.49 & 76.88 & 81.80 & 88.30 & 84.90 & 91.27 \\
  & Cascade-MiniLM & 90.07 & 94.14 & 79.73 & 86.85 & 85.10 & 91.41 \\
  & Cascade-SmolLM & 84.76 & 90.98 & 40.28 & 45.00 & 51.20 & 64.43 \\
  \midrule
  \multirow{4}{*}{$0.50$}
  & Full Pipeline  & \textbf{86.78} & \textbf{91.96} & \textbf{85.54} & \textbf{89.71} & \textbf{78.98} & \textbf{88.23} \\
  & Without Gate   & 67.90 & 70.24 & 80.28 & 82.79 & 77.15 & 82.88 \\
  & Cascade-MiniLM & 85.89 & 90.67 & 77.46 & 79.41 & 77.11 & 84.63 \\
  & Cascade-SmolLM & 73.58 & 78.69 & 50.10 & 55.68 & 65.83 & 77.61 \\
  \midrule
  \multirow{4}{*}{$0.75$}
  & Full Pipeline  & \textbf{79.83} & \textbf{87.10} & \textbf{81.32} & \textbf{75.57} & \textbf{79.83} & \textbf{86.49} \\
  & Without Gate   & 73.40 & 66.50 & 77.36 & 65.83 & 76.50 & 84.20 \\
  & Cascade-MiniLM & 79.10 & 81.00 & 75.91 & 62.50 & 77.23 & 83.69 \\
  & Cascade-SmolLM & 74.31 & 71.79 & 58.91 & 28.72 & 52.48 & 55.41 \\
  \bottomrule
  \end{tabular*}
  \caption{Ablation results under different KIR settings.}
  \label{tab:ablation_all_kir}
\end{table*}

Ablation results are presented in table~\ref{tab:ablation_all_kir}. The full pipeline method outperforms all ablated variants. Compared to these variants, it achieves $2.29\%\sim21.72\%$ improvement on OOS F1 scores across all settings. The results show that using different models is better for OOS detection.

As shown in table~\ref{tab:ablation_all_kir}, OOS F1 scores decrease if removing the front-end gate. That means the post-hoc confidence thresholding is insufficient under complex semantic overlap. A reason is that the unknown queries are easily projected into known-intent regions without front-end rejection. Moreover, the homogeneous variants consistently underperform the full pipeline. That means, if using the same representation model, the scores are suboptimal for cascaded OOS detection. Compared to Cascade-MiniLM, Cascade-SmolLM consistently presents lower OOS F1 scores. A possible reason is that the large generative representations badly perform on the localized geometric rejection, but the MiniLM embeddings preserve more stable local semantic structures for boundary estimation, while the large models introduce additional semantic interference in dense semantic spaces.

\section{Conclusion}

This work proposes a multi-cluster boundary learning method for OOS intent detection via a cascade MiniLM workflow. That workflow separates the OOS intent detection from the multi-class classification of known intents, and views it as a one-class classification problem. Then, multi-cluster boundary learning is conducted to learn the boundary of MiniLM embeddings from the input user utterances. That boundary separates the embedding domain of OOS intents and known intents. Extensive experiments are conducted on public datasets. The results show that the proposed method achieve the stete-of-the-art OOS intent detection performance. Ablation studies also show that the MiniLM settings are reasonable at present. Moreover, the performance of the method is not always the best one on the task of known intent detection, but the cascade workflow completely decouples the two detection process for OOS and known intents. Thus, the result does not influence the aforementioned conclusion. And moreover, the lightweight models have outstanding advantages for real-world applications compared to the large models. The more advanced MiniLMs perhaps improve the detection performance by using this framework in the future.       




\section*{Limitations}
Although the proposed method achieves state-of-the-art performance in OOS detection, its performance on known intents is not state-of-the-art. Specifically, the strict boundary filtering occasionally misclassifies marginal in-domain samples as OOS. This over-rejection slightly compromises Known F1 scores. Additionally, the method relies on the quality of pretrained semantic embeddings, which may become unstable if semantically similar intents are not sufficiently separated. Finally, the evaluation is restricted to monolingual English and single-intent datasets, leaving multilingual and multi-intent extensions for future work.

\section*{Ethics Statement}
This work aims to improve the safety and robustness of human-AI interaction in dialogue systems by rejecting OOS queries. However, automated intent detection should not replace human judgment in high-stakes moderation decisions. Deployment in these scenarios should include appropriate human oversight. Furthermore, since user utterances may contain sensitive personal data, standard privacy and data protection regulations must be strictly maintained.


\bibliography{custom}

\end{document}
